\section{Inference at evaluation time}
\label{sec:method-inference}

At evaluation time the trained predictor maps an input to a
MAP labeling by maximising the joint score
$f(x, y)$ of Eq.~\eqref{eq:method-mn-score} over
$y \in \mathcal{Y}^{|V|}$. The choice of decoder is dictated
by the graph structure: chains admit an exact polynomial-time
decoder, while cyclic graphs require an approximate method.
This section names the concrete decoders used and points to
where in the theoretical background each is derived. The
decoder is selected by the YAML field
\texttt{training.inference.eval} and dispatched by
\texttt{mnlearn.inference}.

\subsection{Chains: Viterbi}
\label{subsec:method-inference-viterbi}

On chain graphs the inner $\arg\max$ is solved exactly in
$O(T K^{2})$ time by the Viterbi algorithm of
Section~\ref{subsec:viterbi}, implemented in
\texttt{mnlearn.inference.viterbi\_decode} and selected by
\texttt{inference.eval: viterbi} in the experiment YAML.

\subsection{Cyclic graphs: Belief Propagation}
\label{subsec:method-inference-cyclic}

On the Sudoku constraint graph, MAP is NP-hard
(Section~\ref{sec:map-inference}). The decoder used in this
thesis is the loopy max-sum Belief Propagation algorithm of
Section~\ref{subsec:belief-propagation}, implemented in
\texttt{mnlearn.inference.bp\_decode} with messages re-normalised
by subtracting the per-edge maximum after each update for
numerical stability. The number of message-passing iterations is
a YAML knob (\texttt{inference.params.bp\_iters}, default $50$);
the decoder is selected by \texttt{inference.eval: bp}.